
\chapter{Khi Thầy Cô Cũng Vô Lý}
\markboth{Khi Thầy Cô Cũng Vô Lý}{Khi Thầy Cô Cũng Vô Lý}

\section{``Đi Học Muộn Là Lỗi Em, Kẹt Xe Không Phải Lý Do''}
\begin{truyen}{``Đi Học Muộn Là Lỗi Em, Kẹt Xe Không Phải Lý Do''}{``Late Is Your Fault. Traffic Is Not an Excuse''}
\chuhoa{H}{ọc sinh đến lớp trễ, thở hổn hển:} ``Thưa thầy, em kẹt xe. Khu vực trước trường hôm nay tắc từ đầu đường ạ.''
\textit{(Student arrives panting: ``I was stuck in traffic, teacher. The road near school was blocked from the start.'')}

\teacherpair{Đi sớm hơn thì không kẹt.}{``Leave earlier and you would not be stuck.''}

\studentpair{Nhưng thầy ơi, hôm nay tai nạn bất ngờ mà.}{``But teacher, there was a sudden accident today.''}

\teacherpair{Thầy không muốn nghe lý do. Đứng ngoài cửa.}{``I do not want to hear reasons. Stand outside.''}

Học sinh đứng ngoài cửa trong một tiết quan trọng. Không phải vì nó không cố. Mà vì giáo viên không muốn phân biệt giữa lý do chính đáng và bào chữa.
\textit{(The student stands outside during an important lesson. Not because they failed to try. But because the teacher did not want to distinguish between a valid reason and an excuse.)}

\begin{baihoc}
Có lý do thật và có bào chữa. Kẹt xe đột xuất vì tai nạn không giống với dậy muộn vì xem phim khuya. Một giáo viên khôn ngoan biết hỏi thêm một câu trước khi phán xét.
\textit{(There are real reasons and there are excuses. An unexpected traffic accident is not the same as oversleeping after watching a late movie. A wise teacher asks one more question before judging.)}
\end{baihoc}
\end{truyen}

\section{Điểm Trừ Vì Chữ Xấu}
\begin{truyen}{Điểm Trừ Vì Chữ Xấu}{Points Deducted for Bad Handwriting}
\chuhoa{H}{ọc sinh:} ``Bài em đúng hết nhưng thầy trừ điểm chữ xấu. Mà thầy không báo trước tiêu chí đó.''
\textit{(Student: ``My answers were all correct but you deducted points for bad handwriting. You never mentioned that criterion before.'')}

\teacherpair{Chữ đẹp là thái độ. Thái độ tính điểm được.}{``Good handwriting reflects attitude. Attitude counts toward grades.''}

\studentpair{Thưa thầy, nhưng lần sau thầy có thể nói trước được không để tụi em biết mà chuẩn bị?}{``Could you tell us in advance next time so we can prepare?''}

\teacherpair{Đáng lý phải biết từ đầu rồi. Không cần phải nhắc.}{``You should have known from the start. There was no need for a reminder.''}

Học sinh ghi lại để hỏi ban giám hiệu. Không phải để kiện. Chỉ để biết mình có quyền được thông báo tiêu chí hay không.
\textit{(The student makes a note to ask the principal's office. Not to file a complaint. Just to find out whether they have a right to be informed of grading criteria.)}

\begin{baihoc}
Tiêu chí bất ngờ là tiêu chí bất công. Giáo viên có mọi quyền đánh giá thái độ. Nhưng tiêu chí đó phải được thông báo trước và áp dụng nhất quán cho cả lớp.
\textit{(Surprise criteria are unfair criteria. Teachers have every right to assess attitude. But that criterion must be communicated in advance and applied consistently to the entire class.)}
\end{baihoc}
\end{truyen}

\section{``Tôi Không Chịu Trách Nhiệm Chuyện Nhà Em''}
\begin{truyen}{``Tôi Không Chịu Trách Nhiệm Chuyện Nhà Em''}{``That Is a Home Problem, Not My Responsibility''}
\chuhoa{H}{ọc sinh đến tìm thầy chủ nhiệm:} ``Thưa thầy, dạo này em không tập trung nổi vì ở nhà có chuyện. Thầy có thể thông cảm cho em một thời gian được không?''
\textit{(Student comes to the homeroom teacher: ``I have not been able to focus because of something at home. Can you be patient with me for a while?'')}

\teacherpair{Chuyện nhà là chuyện nhà. Vào lớp là chuyện học. Tôi không can thiệp chuyện nhà.}{``Home is home. School is school. I do not interfere with home matters.''}

Đúng. Thầy không thể giải quyết chuyện nhà học sinh. Nhưng ``thông cảm'' không phải ``giải quyết''. Đứa trẻ không xin được sửa nhà. Nó chỉ xin được nhìn nhận.
\textit{(True. The teacher cannot fix a student's home situation. But ``understanding'' is not ``solving.'' The student did not ask for the problem to be fixed. They only asked to be seen.)}

\begin{baihoc}
Giáo viên không phải nhà tâm lý học. Nhưng giáo viên là người có mặt trong đời học sinh nhiều giờ mỗi tuần. Nói ``không phải việc tôi'' khi học sinh dám mở miệng tâm sự là một cách đóng cửa sẽ khó mở lại.
\textit{(Teachers are not psychologists. But they are present in students' lives for many hours every week. Saying ``not my job'' when a student dares to open up is a way of closing a door that will be very hard to reopen.)}
\end{baihoc}
\end{truyen}

\section{Cáu Vì Ngoài Trời Nóng, Đổ Lên Học Sinh}
\begin{truyen}{Cáu Vì Ngoài Trời Nóng, Đổ Lên Học Sinh}{Irritated by the Heat, Taking It Out on Students}
\chuhoa{T}{rời nóng 38 độ.} Điều hoà lớp hỏng ba ngày. Cả thầy lẫn trò đang chảy mồ hôi. Một học sinh đùa một câu nhỏ trong lớp. Giáo viên đang trong trạng thái cực kỳ khó chịu.
\textit{(The temperature is 38 degrees. The classroom air conditioner has been broken for three days. Both teacher and students are sweating. A student cracks a small joke. The teacher is in an extremely irritated state.)}

\teacherpair{Tưởng mình buồn cười lắm hả? Ra ngoài ngồi đi. Chọc phá cả buổi.}{``Think you are funny? Go sit outside. You have been disrupting the whole session.''}

\studentpair{Nhưng thầy ơi, em chỉ nói một câu nhỏ thôi.}{``But teacher, I only said one small thing.''}

\teacherpair{Một câu là đủ rồi. Ra ngoài.}{``One sentence is enough. Outside.''}

Ba học sinh cùng bàn nhìn nhau. Cùng nghĩ một điều: Trời nóng chứ không phải em nó. Nhưng không ai dám nói.
\textit{(Three students at the same table look at each other. All thinking the same thing: It is the heat, not the student. But no one dares to say it.)}

\begin{baihoc}
Giáo viên cũng là người. Cũng mệt, cũng nóng, cũng cáu. Điều đó hoàn toàn bình thường.\\\
Điều không bình thường là không nhận ra mình đang đổ cảm xúc của mình lên người không liên quan.
\textit{(Teachers are human. They also feel tired, hot, and irritable. That is entirely normal. What is not normal is failing to recognize when you are pouring your emotions onto someone uninvolved.)}
\end{baihoc}
\end{truyen}

\section{``Những Đứa Không Học Thêm Thầy Đừng Trông Chờ Điểm Cao''}
\begin{truyen}{``Những Đứa Không Học Thêm Thầy Đừng Trông Chờ Điểm Cao''}{``Don't Expect High Marks If You Don't Attend My Extra Classes''}
\chuhoa{P}{hụ huynh phát hiện} điểm kiểm tra con mình bắt đầu tụt sau khi không đăng ký học thêm với giáo viên bộ môn. Hỏi con, đứa học sinh kể lại câu giáo viên nói đầu năm:
\textit{(A parent notices the child's test scores dropping after declining the teacher's private tutoring classes. Asking the child, the student recalls what the teacher said at the beginning of the year:)}

\textit{``Những đứa không học thêm thầy thì thầy không đặt kỳ vọng cao. Vì thầy không biết năng lực của tụi nó đến đâu.''}
\textit{(``Students who do not take my extra classes, I do not hold high expectations for. Because I do not know their level.'')}

Phụ huynh ngồi lặng một hồi. Rồi viết thư gửi hiệu trưởng.
\textit{(The parent sits in silence for a moment. Then writes a letter to the principal.)}

\begin{baihoc}
Dạy thêm là nhu cầu có thật. Học sinh cần thêm hỗ trợ. Thầy cô cũng cần thêm thu nhập. Điều đó có thể được tôn trọng. Nhưng biến lớp học thêm thành điều kiện để được đối xử công bằng trên lớp chính là một đường ranh mà không có giáo viên nào nên bước qua.
\textit{(Extra tutoring has a genuine purpose. Students need support. Teachers also need income. That can be respected. But making tutoring classes a condition for fair treatment in the main classroom is a line no teacher should cross.)}
\end{baihoc}
\end{truyen}

\section{Nhớ Lỗi Cũ, Không Cho Cơ Hội Mới}
\begin{truyen}{Nhớ Lỗi Cũ, Không Cho Cơ Hội Mới}{Memorizing Old Faults, Refusing New Chances}
\chuhoa{H}{ọc sinh đầu học kỳ hai} đã cố gắng thay đổi: nộp bài đúng hạn, không nói chuyện, giơ tay phát biểu.

\textit{(At the start of the second semester, a student makes a deliberate effort to change: submitting on time, staying quiet, raising their hand to answer.)}

\teacherpair{Thầy biết em sẽ lại như cũ thôi. Đợt trước em cũng hay lắm rồi quên.}{``I know you will go back to your old ways. Last time you were also well-behaved for a while and then forgot.''}

Học sinh gật đầu. Cố lịch sự. Nhưng trong lòng, có gì đó vừa xì hơi.
\textit{(The student nods. Tries to stay polite. But inside, something just deflated.)}

\begin{baihoc}
Sự thay đổi cần được nhìn nhận. Không phải khen ngợi quá mức, nhưng ít nhất là một cơ hội thật sự. Giáo viên ghi nhớ lỗi học sinh lâu hơn ghi nhận sự cố gắng của chúng thì vô tình trở thành rào cản của chính sự thay đổi mà họ mong muốn.
\textit{(Change needs to be acknowledged. Not excessively praised, but at least given a real chance. Teachers who remember a student's faults longer than they notice the student's effort can inadvertently become the biggest obstacle to the very change they want to see.)}
\end{baihoc}
\end{truyen}

\begin{baihoc}
``Khi thầy cô cũng vô lý'' không phải chương viết để bêu xấu giáo viên. Đó là chương được viết vì phòng học cần sự trung thực từ cả hai phía.
\textit{(``When teachers are also unreasonable'' is not a chapter written to shame teachers. It is written because the classroom needs honesty from both sides.)}

Giáo viên có áp lực rất thật, có nỗi mệt rất thật, và có giới hạn rất thật. Nhưng những điều đó không tự động cho phép họ miễn trừ khỏi tiêu chuẩn đúng đắn.
\textit{(Teachers face very real pressure, very real fatigue, and very real limits. But none of those automatically exempt them from being held to a reasonable standard.)}
\end{baihoc}

\begin{ghinhoanh}
\textbf{Short English to remember:}

Fairness has no uniform.

Students remember how teachers made them feel, long after they forget what was taught.
\end{ghinhoanh}

\begin{tuhocnhanh}
1. Điều gì khiến em cảm thấy bị đối xử không công bằng nhất trong lớp học? \textit{(What made you feel most unfairly treated in a classroom?)}

2. Làm sao phân biệt ``thầy cô đang vô lý'' với ``em chưa đủ trưởng thành để hiểu''? \textit{(How do you distinguish ``the teacher is being unreasonable'' from ``I am not yet mature enough to understand''?)}
\end{tuhocnhanh}

\ngancach
